% ============================================
% Johns Hopkins Letter Template
% Assistant Professor Cover Letter – Duquesne University
% ============================================

\documentclass[11pt,letterpaper]{letter}

\usepackage{graphicx}
\usepackage{eso-pic}
\usepackage{geometry}
\usepackage{microtype}
\usepackage[hidelinks]{hyperref}

% ----- Page Setup -----
\geometry{left=25mm,right=25mm,top=25mm,bottom=25mm}

% ----- Logo Placement (foreground, first page only) -----
\newcommand{\PutJHULogoFG}[1]{%
  \AddToShipoutPictureFG*{%
    \AtPageUpperLeft{%
      \hspace{10mm}%
      \raisebox{-38mm}{%
        \includegraphics[width=6cm]{#1}%
      }%
    }%
  }%
}

% ----- Sender Information -----
\signature{Decory Edwards}
\address{Department of Economics \\ Johns Hopkins University \\ Baltimore, MD 21218 \\ \texttt{dedwar65@jh.edu}}

\begin{document}
\PutJHULogoFG{Images/jhulogo.png}

\begin{letter}{Search Committee \\
Department of Economics and Finance \\
Palumbo-Donahue School of Business \\
Duquesne University \\
Pittsburgh, PA \\ USA}

\opening{Dear Members of the Search Committee,}

I am writing to express my interest in the tenure-track Assistant Professor of Economics position in the Palumbo-Donahue School of Business at Duquesne University, beginning in August 2026. I am a Ph.D. candidate in Economics at Johns Hopkins University, advised by Professors Christopher D. Carroll, Nicholas W. Papageorge, and Stelios Fourakis, and I expect to receive my degree in May 2026. My research fields are macroeconomics, wealth and income dynamics, and computational economics.

My research agenda focuses on understanding the determinants of wealth inequality and the mechanisms through which heterogeneity in returns to wealth shapes aggregate outcomes. In my job-market paper, I estimate the degree of heterogeneity in individual portfolio returns using micro-data from the Survey of Consumer Finances and simulate a structural model in which households face idiosyncratic return risk. I show that allowing for persistent heterogeneity in returns substantially improves the model’s ability to match the empirical distribution of wealth, offering a tractable way to connect micro-level return dispersion to macroeconomic inequality.

In a second project, I use the Health and Retirement Study to examine the relationship between interpersonal trust and portfolio performance. Building on the literature that documents a hump-shaped relationship between trust and income, I show that a similar relationship holds for individual returns to wealth. This project highlights a behavioral channel through which social attitudes shape saving and investment outcomes and provides new evidence on the persistence of return heterogeneity, which directly informs the calibration of my structural framework.

I am particularly interested in Duquesne University because of the Palumbo-Donahue School of Business’s emphasis on integrating rigorous economic analysis with applied questions relevant to finance, sustainability, and ethical business practice. My work on heterogeneous portfolio returns and household risk exposure connects naturally to the School’s strengths in financial economics and to initiatives such as the Investment Strategy Institute, where empirical evidence on risk, returns, and long-run outcomes is central. I am also drawn to Duquesne’s research priorities in sustainable business practices and business ethics. Understanding how inequality, risk, and market structure shape household and firm behavior is directly relevant to these areas, and I would welcome opportunities for interdisciplinary collaboration with colleagues working on entrepreneurship, strategy, and socially responsible markets.

\textbf{Teaching.} My teaching experience centers on undergraduate macroeconomics and an upper-level macro policy seminar, where I have led weekly discussion sections and held regular office hours for classes ranging from 16 to 40 students. My approach is student-centered. I focus on helping students convince themselves that they have moved from confusion to understanding, using structured problem-solving and coaching in office hours. At Duquesne, I am prepared to teach undergraduate and graduate courses in macroeconomics, principles-level economics, and quantitative methods, and I am comfortable contributing to statistics or applied economics courses that serve business students.

I believe I would be a strong fit for the Department of Economics and Finance because my research and teaching align with Duquesne’s teacher-scholar model. My work combines careful empirical analysis, structural modeling, and policy-relevant questions, while my teaching emphasizes clarity, engagement, and real-world application. I value collegial service and student mentorship, and I would be eager to contribute to the School’s academic programs, research initiatives, and broader mission. I also respect and appreciate Duquesne’s Spiritan Catholic identity, particularly its commitment to service, ethical reflection, and concern for social outcomes, values that align naturally with my research focus on inequality and economic well-being.

I would be honored to join Duquesne University and contribute to the Palumbo-Donahue School of Business through teaching, research, and service. Thank you for your consideration. I welcome the opportunity to discuss my candidacy with you.

\closing{Sincerely,}

\end{letter}
\end{document}
